%==============================================================================
% Mahasweta Bhattacharya, PhD — Resume (BigTech AI/ML Scientist version)
% Single-column, ATS-clean. Compile with pdflatex or xelatex on Overleaf.
%==============================================================================
\documentclass[11pt,letterpaper]{article}

\usepackage[margin=0.65in]{geometry}
\usepackage[T1]{fontenc}
\usepackage[utf8]{inputenc}
\usepackage{lmodern}
\usepackage[hidelinks]{hyperref}
\usepackage{enumitem}
\usepackage{titlesec}
\usepackage{xcolor}
\usepackage{microtype}

\definecolor{accent}{HTML}{1F4E79}

% Compact section headers, ATS-friendly (text remains plain)
\titleformat{\section}{\large\bfseries\color{accent}}{}{0pt}{}[\vspace{-4pt}\hrule height 0.4pt\vspace{2pt}]
\titlespacing*{\section}{0pt}{8pt}{4pt}

\titleformat{\subsection}{\normalsize\bfseries}{}{0pt}{}
\titlespacing*{\subsection}{0pt}{6pt}{2pt}

\setlist[itemize]{leftmargin=1.1em,topsep=2pt,parsep=0pt,itemsep=2pt}
\setlength{\parindent}{0pt}
\pagenumbering{gobble}

\newcommand{\role}[4]{%
  \noindent\textbf{#1} \hfill \textit{#3}\\[1pt]
  \textit{#2} \hfill \textit{#4}\par\vspace{2pt}}

\begin{document}

%------------------- HEADER -------------------
\begin{center}
  {\LARGE\bfseries Mahasweta Bhattacharya, PhD}\\[3pt]
  {AI/ML Research Scientist \textbar{} Drug Discovery \textbar{} Foundation Models \textbar{} Agentic AI}\\[4pt]
  Malden, MA \textbar{} \href{mailto:b.mahasweta24@gmail.com}{b.mahasweta24@gmail.com} \textbar{} 609-906-1583\\
  \href{https://www.linkedin.com/in/mahasweta-bhattacharya/}{linkedin.com/in/mahasweta-bhattacharya} \textbar{}
  \href{https://github.com/mahaswetabhattacharya24}{github.com/mahaswetabhattacharya24} \textbar{}
  \href{https://mahaswetabhattacharya24.github.io}{mahaswetabhattacharya24.github.io}
\end{center}

%------------------- SUMMARY -------------------
\section*{Summary}
AI/ML Research Scientist with 3+ years of industry experience and a PhD in Biomedical Engineering, designing
foundation models, agentic LLM pipelines, and graph neural networks for biological discovery at scale.
At Sanofi, built multimodal ML systems achieving 4$\times$ improvement in causal target recall, reduced
manual target review by $>$50\% via multi-agent orchestration, and contributed to advancing 7 novel targets
into preclinical evaluation across 30+ therapeutic programs. Prior work includes GPU-accelerated proton
dose engines (0.28s vs 4.68s Monte Carlo) and deep reinforcement learning for radiotherapy. Published
across peer-reviewed journals; presenter at international AI and biomedical conferences.

%------------------- SKILLS -------------------
\section*{Skills}

\textbf{Machine Learning \& AI:} Foundation Models, Multimodal Learning, Large Language Models (LLMs),
Fine-tuning (LoRA / PEFT), Retrieval-Augmented Generation (RAG), Agentic AI, Multi-Agent Orchestration,
Graph Neural Networks (GNNs), Causal Inference, Deep Reinforcement Learning (RLHF / VMAT),
Transformer Architectures, Representation Learning, Variational Autoencoders, Bayesian Inference,
Uncertainty Quantification.\\[3pt]
\textbf{Frameworks \& Tools:} PyTorch, TensorFlow, JAX, Hugging Face Transformers, LangChain,
scikit-learn, CUDA / GPU computing, Docker, Kubernetes, Slurm HPC, Weights \& Biases, MLflow.\\[3pt]
\textbf{Cloud \& Data:} AWS (S3, EC2, Batch, SageMaker), Snowflake, Neo4j, MySQL, FAISS / vector search,
S3 data lakes, ETL pipelines, data harmonization \& QC.\\[3pt]
\textbf{Languages:} Python (expert), R, MATLAB, SQL, Bash / Shell, C / C++.\\[3pt]
\textbf{Bioinformatics \& Drug Discovery:} Bulk \& single-cell RNA-seq, GWAS \& multi-omics integration,
target credentialing, indication discovery, pharmacodynamics modeling, biological knowledge graphs,
embedding similarity search, scientific evidence synthesis.

%------------------- EXPERIENCE -------------------
\section*{Experience}

\role{Senior Scientist --- AI / ML Research}{Sanofi --- Cambridge, MA}{Sept 2023 -- Present}{}
\begin{itemize}
  \item Designed and deployed a production multi-agent LLM pipeline orchestrating Claude-based agents
        that autonomously aggregate genetics, transcriptomics, and clinical literature into structured
        target credentialing reports; reduced expert manual review time by $>$50\% across 30+ therapeutic
        programs and operationalized LLM-driven reasoning in a regulated scientific workflow.
  \item Architected a multimodal foundation model integrating GWAS embeddings, bulk and single-cell
        transcriptomics, and clinical-trial outcome embeddings into a unified representation space;
        achieved 4$\times$ improvement in causal target recall over genetics-only baselines and established
        a transferable backbone for cross-disease generalization.
  \item Founded a graph neural network platform for bispecific antibody discovery integrating
        protein-protein interaction graphs, synergy metrics, and LLM-guided evidence retrieval;
        generated 5 novel bispecific target-pair candidates supporting pipeline expansion decisions.
  \item Built a scalable indication-discovery engine scoring 232 immune indications in 3 weeks and
        scaling to 17,000+ phenotypes via parallelized Slurm/HPC jobs; enabled systematic whitespace
        identification across the immunology portfolio.
  \item Developed an explainable AI target-discovery engine producing 90+ ranked hypotheses and
        contributing to 7 novel targets entering preclinical evaluation; integrated causal scoring,
        embedding similarity, and LLM-augmented evidence synthesis.
  \item Co-led an automated multimodal target-credentialing platform supporting 30+ therapeutic
        programs and contributing to 3 preclinical nominations; introduced modules for causal inference,
        uncertainty quantification, and prospective validation.
  \item Led ML-driven transcriptomic pharmacodynamics modeling for Hidradenitis Suppurativa,
        differentiating oral vs.\ injectable therapy mechanisms and enabling preclinical advancement
        of the oral candidate.
\end{itemize}

\role{Postdoctoral Fellow --- Computational AI \& Medical Physics}{Johns Hopkins University School of Medicine --- Baltimore, MD}{Jan 2023 -- Sept 2023}{}
\begin{itemize}
  \item Built a GPU-accelerated proton pencil beam dose engine in Python / CUDA achieving 0.28s mean
        patient-plan computation vs.\ 4.68s Monte Carlo (96\% 3D-gamma agreement at 2\%/2mm); released
        as an open-source Python package and published in \textit{Journal of Applied Clinical Medical
        Physics} (2025).
  \item Architected an end-to-end beam-modeling and validation pipeline from 98 pristine Bragg-peak
        energies, benchmarking against measurements and heterogeneous digital phantoms; produced
        uncertainty profiles highlighting limits in heterogeneous regions.
  \item Developed a deep reinforcement learning agent for VMAT machine-parameter optimization;
        produced deliverable prostate plans in 3.3s and reduced mean rectum dose vs.\ clinical plans
        (17.4 vs.\ 21.0 Gy, $p=0.024$) on a 15-patient external cohort. Published in
        \textit{Medical Physics} (2024).
  \item Integrated RL-generated optimization with a commercial treatment-planning system into a
        hybrid automated + human-in-the-loop workflow, demonstrating production-grade feasibility.
\end{itemize}

\role{Research Scientist (PhD) --- Computational Neuroscience \& AI}{State University of New York at Buffalo --- Buffalo, NY}{2017 -- 2023}{}
\begin{itemize}
  \item Reconstructed session-wise functional connectomes from longitudinal two-photon calcium imaging
        during motor-skill learning; identified a biphasic rewiring trajectory and discovered stable
        L2/3 hub assemblies encoding movement.
  \item Developed FARCI, an open-source MATLAB toolbox for fast connectome inference using OASIS spike
        deconvolution and partial correlations; outperformed state-of-the-art on NCC and NAOMi
        benchmarks in accuracy and scalability.
  \item Demonstrated low-frequency frontal HbO--EEG coupling (0.07--0.13 Hz) as a consciousness
        biomarker in acute brain injury; AMICA embeddings + k-NN classifier achieved $>$90\% accuracy
        distinguishing conscious vs.\ unresponsive patients and $>$99\% accuracy predicting failure to
        recover consciousness.
  \item Integrated portable fNIRS, EEG, and electric-field--informed GLMs to model cerebellar tDCS
        response in stroke; uncovered hemoglobin--EEG coupling signatures predictive of responders.
  \item Built Monte Carlo photon-transport and bioheat models for transcranial near-infrared
        stimulation (810 nm); supported safety and CCO-mediated photobiomodulation as the primary
        mechanism --- published in \textit{Brain Sciences} (2019).
\end{itemize}

%------------------- EDUCATION -------------------
\section*{Education}
\textbf{PhD, Biomedical Engineering} \hfill 2017 -- 2023\\
University at Buffalo, State University of New York --- Buffalo, NY\\
{\small Dissertation: neural signal processing, functional connectomics, brain--computer interfaces, and computational photobiomodulation modeling.}\\[3pt]
\textbf{MS, Electrical Engineering} \hfill 2015 -- 2017\\
University at Buffalo, State University of New York --- Buffalo, NY\\[3pt]
\textbf{BTech, Electronics and Communication Engineering} \hfill 2010 -- 2014\\
West Bengal University of Technology --- Kolkata, India

%------------------- PUBLICATIONS -------------------
\section*{Selected Publications}
\begin{enumerate}[leftmargin=1.4em,itemsep=2pt,topsep=2pt]
  \item \textbf{Bhattacharya M}, Reamy C, Li H, Lee J, Hrinivich WT.
        ``A Python package for fast GPU-based proton pencil beam dose calculation.''
        \textit{Journal of Applied Clinical Medical Physics}, 2025.
  \item Hrinivich WT, \textbf{Bhattacharya M}, et al. ``Clinical VMAT machine parameter optimization
        for localized prostate cancer using deep reinforcement learning.''
        \textit{Medical Physics}, 51(6): 3972--3984, 2024.
  \item Meamardoost S, \textbf{Bhattacharya M}, et al. ``Rewiring dynamics of functional connectomes
        during motor-skill learning.'' \textit{Data Science in Science}, 2(1): 2260431, 2023.
  \item Meamardoost S, \textbf{Bhattacharya M}, et al. ``FARCI: Fast and robust connectome inference.''
        \textit{Brain Sciences}, 11(12): 1556, 2021.
  \item \textbf{Bhattacharya M}, Dutta A. ``Computational modeling of the photon transport, tissue
        heating, and cytochrome c oxidase absorption during transcranial near-infrared stimulation.''
        \textit{Brain Sciences}, 9(8): 179, 2019.
\end{enumerate}
{\small Additional publications in \textit{Neurocritical Care} (2021), \textit{The Cerebellum} (2021),
and \textit{Scientific Reports} (2020). Full list at
\href{https://mahaswetabhattacharya24.github.io/publications/}{mahaswetabhattacharya24.github.io/publications}.}

%------------------- TALKS -------------------
\section*{Selected Talks \& Presentations}
\begin{itemize}
  \item \textbf{HubXChange: AI in Drug Discovery}, Boston, MA --- Speaker / Panelist (2025)
  \item \textbf{ISMB / ECCB}, Liverpool, UK --- Co-author, \textit{MultiTIE: Sanofi's Multi-Specific Target Immune Engine} (2025)
  \item \textbf{FOCIS}, Boston, MA --- Co-author, \textit{HSID: Integrative Target Discovery for Hidradenitis Suppurativa} (2025)
  \item \textbf{AAPM Annual Meeting}, Houston, TX --- Presenter, \textit{GPU-Based Proton Pencil Beam Dose Calculation} (2023)
  \item \textbf{Society for Neuroscience Annual Meeting}, San Diego, CA --- Presenter, \textit{Stability of Motor Cortex Decoders during Learning} (2022)
\end{itemize}

%------------------- HONORS -------------------
\section*{Honors, Service \& Community}
\textbf{Awards.} Winning Solution, Biomedical Knowledge Graph Hackathon, BioLabs Heidelberg \& Sanofi (2025);
Finalist, 3 Minute Thesis Competition (2022).\\[3pt]
\textbf{Memberships.} IEEE Senior Member (elevated 2025); The AI Consortium (2025--).\\[3pt]
\textbf{Peer Review.} Applied Sciences, BioMedInformatics, Engineering Applications of Artificial
Intelligence (Elsevier), Heliyon, International Journal of Molecular Sciences, Mathematics, Symmetry,
Genes.\\[3pt]
\textbf{Mentorship \& Teaching.} MIT UPOP Mentor Panelist (2024--); Hopkins Biotech Network Mentor Match
(2024); Course Instructor, Biomedical Circuits and Signals \& Biosignals Lab, SUNY Buffalo (2017--2019).

\end{document}
